\section*{Hoofdstuk \ref{ch:b1:organic-redox} --- Oxidatie en reductie in de organische chemie}

\begin{solution}{exo:b1:organic-redox:1}
Ethanol: \ce{CH3} $-3$, \ce{CH2OH} $-1$. Ethanal: \ce{CH3} $-3$, CHO $+1$.
Ethaanzuur: \ce{CH3} $-3$, COOH $+3$. Propanon: \ce{CH3} $-3$ (twee keer),
C=O $+2$. Methaanzuur: $+2$.
\end{solution}

\begin{solution}{exo:b1:organic-redox:2}
Propaan-1-ol: propanal, daarna propaanzuur. Propaan-2-ol: propanon.
2-Methylpropaan-2-ol: geen reactie (de oranje kleur blijft).
\end{solution}

\begin{solution}{exo:b1:organic-redox:3}
\ce{CH3COCH3 + 2H+ + 2e- -> CH3CH(OH)CH3};
\ce{CH3CHO + 2H+ + 2e- -> CH3CH2OH}.
\end{solution}

\begin{solution}{exo:b1:organic-redox:4}
\ce{NaBH4}: butaan-1-ol, butaan-2-ol, geen reactie, geen reactie (het zuur
wordt alleen gedeprotoneerd). \ce{LiAlH4}: butaan-1-ol, butaan-2-ol,
butaan-1-ol en ethanol, butaan-1-ol.
\end{solution}

\begin{solution}{exo:b1:organic-redox:5}
\ce{5CH3CH(OH)CH3 + 2MnO4- + 6H+ -> 5CH3COCH3 + 2Mn^2+ + 8H2O}.
$n = 1.20/60.0 = \qty{0.0200}{mol}$; permanganaat $\frac25 \times 0.0200 =
\qty{8.0e-3}{mol}$, \qty{400}{mL} van de oplossing van \qty{0.020}{mol/L}.
\end{solution}

\begin{solution}{exo:b1:organic-redox:6}
Verwarm de alcohol met de \omterm{def:b1:nernst:couple}{oxidator} in een kolf die voor destillatie
opgesteld is, bij een temperatuur tussen de twee kookpunten: propanal
(\qty{48}{\celsius}) destilleert over naarmate het ontstaat en ontsnapt
aan verdere oxidatie; propaan-1-ol (\qty{97}{\celsius}) blijft in de
kolf.
\end{solution}

\begin{solution}{exo:b1:organic-redox:7}
Een B--H-paar van \ce{BH4-} valt het carbonylkoolstofatoom aan, terwijl
het $\pi$-paar van C=O naar de zuurstof gaat; het \omterm{def:b1:alcohol-activation:alkoxide}{alkoxide} neemt een
proton van ethanol op. Het waterstofatoom op het koolstofatoom van het
product (\ce{CH3CH2OH}, een van de twee op C1) komt van het reagens; dat
op de zuurstof van het oplosmiddel.
\end{solution}

\begin{solution}{exo:b1:organic-redox:8}
Butaan-2-ol is \omterm{def:b1:stereochemistry-in-depth:stereocentre}{chiraal}, maar het vlakke keton wordt even gemakkelijk aan
beide zijden aangevallen: een \omterm{def:b1:stereochemistry-in-depth:enantiomers}{racemisch mengsel}, optisch inactief.
\end{solution}

\begin{solution}{exo:b1:organic-redox:9}
\ce{3CH3CH2OH + 2Cr2O7^2- + 16H+ -> 3CH3COOH + 4Cr^3+ + 11H2O}: oranje
dichromaat (chroom $+$VI) wordt groen chroom(III).
\end{solution}

\begin{solution}{exo:b1:organic-redox:10}
2-Methylbutaan-2-ol: ethanal + \ce{CH3CH2MgBr} geeft butaan-2-ol (een
additie, geen redox); oxidatie tot butanon; \ce{CH3MgBr} op butanon,
daarna opwerking. Eén oxidatie. Pentaan-3-ol: propanal + \ce{CH3CH2MgBr}
in één stap, geen redox.
\end{solution}

\begin{solution}{exo:b1:organic-redox:11}
1. Bescherm het keton als cyclisch \omterm{def:b1:acetals-protection:hemiacetal}{acetaal} (ethaan-1,2-diol, zuur,
Dean--Stark). 2. \ce{LiAlH4} in droge ether reduceert de ester tot de
primaire alcohol (en methanol). 3. Waterig zuur verwijdert het \omterm{def:b1:acetals-protection:hemiacetal}{acetaal}:
5-hydroxypentaan-2-on, \ce{CH3CO(CH2)3OH}.
\end{solution}

\begin{solution}{exo:b1:organic-redox:12}
\ce{RCH2OH}: carbinolkoolstofatoom op $-1$; \ce{RCOOH}: $+3$; vier
elektronen, zoals in \ce{RCH2OH + H2O -> RCOOH + 4H+ + 4e-}. Methaan
($-4$) naar koolstofdioxide ($+4$): acht elektronen,
\ce{CH4 + 2H2O -> CO2 + 8H+ + 8e-}.
\end{solution}

\begin{solution}{pb:b1:organic-redox:1}
\textbf{1.} C1 (draagt OH): 0; de vijf \ce{CH2}: $-2$.
\textbf{2.} C1 (C=O): $+2$; de vijf \ce{CH2}: $-2$.
\textbf{3.} De twee COOH-koolstofatomen: $+3$; de vier \ce{CH2}: $-2$.
\textbf{4.} C1 ($0 \to +3$) en zijn buur C6 ($-2 \to +3$), dat het tweede
carboxylkoolstofatoom wordt; de binding C1--C6 wordt verbroken.
\textbf{5.} $3 + 5 = 8$ elektronen.
\textbf{6.} Salpeterzuur is sterk genoeg om de C--C-binding naast de
carbonylgroep te verbreken, wat mildere \omterm{def:b1:nernst:couple}{oxidatoren} niet doen.
\textbf{7.} \ce{C6H12O + 3H2O -> C6H10O4 + 8H+ + 8e-}.
\textbf{8.} $+$V en $+$I.
\textbf{9.} \ce{2HNO3 + 8H+ + 8e- -> N2O + 5H2O}.
\textbf{10.} Bij het optellen vallen acht elektronen, \ce{8H+} en drie
watermoleculen weg: \ce{C6H12O + 2HNO3 -> C6H10O4 + N2O + 2H2O}.
\textbf{11.} 8.
\textbf{12.} $10^6/146.0 = \qty{6.85e3}{mol}$ zuur; $2 \times 6.85 \times
10^3 \times 63.0 = \qty{8.6e5}{g}$, \qty{0.86}{t} salpeterzuur.
\textbf{13.} \ce{H2O2 + 2H+ + 2e- -> 2H2O}.
\textbf{14.} \ce{C6H12O + 4H2O2 -> C6H10O4 + 5H2O}.
\textbf{15.} $4 \times 6.85 \times 10^3 \times 34.0 = \qty{9.3e5}{g}$,
\qty{0.93}{t}.
\textbf{16.} Water: geen stikstofoxide.
\textbf{17.} Een grote evenwichtsconstante zegt niets over de snelheid;
\ce{H2O2} reageert zonder \omterm{def:b1:elementary-steps:catalyst}{katalysator} traag, net zoals zijn eigen
ontleding.
\textbf{18.} Het reduceert cyclohexanon terug tot cyclohexanol:
\ce{4C6H10O + BH4- -> B(OC6H11)4-}, daarna
\ce{B(OC6H11)4- + 4H2O -> 4C6H11OH + B(OH)4-}.
\textbf{19.} \qty{6.85e3}{mol}.
\textbf{20.} \qty{6.85e3}{mol} \ce{N2O} (één per zuur).
\textbf{21.} $6.85 \times 10^3/0.95 \times 100.0 = \qty{7.2e5}{g}$, \qty{0.72}{t}.
\textbf{22.} $6.85 \times 10^3 \times 44.0 = \qty{3.0e5}{g}$: \textbf{\qty{0.30}{t}
\ce{N2O} per ton adipinezuur}.
\end{solution}
